\chapter{A primer of basic \LaTeX \ commands}\label{app:latex-primer}

An example figure looks like the one below in Figure \ref{fig:example-figure}.

\begin{figure}[ht]
    \centering
    \includegraphics[width=0.75\textwidth]{example-image-a}
    \caption{Here's an example figure}
    \label{fig:example-figure}
\end{figure}

For an example footnote, use this syntax\footnote{\href{https://www.cst.cam.ac.uk/}{The CST Department Website}}.

For bullet points:
\begin{itemize}
    \item First item.
    \item Second item.
\end{itemize}

For enumerated lists:
\begin{enumerate}
    \item First item.
    \item Second item.
\end{enumerate}

See Table \ref{tab:example-table} for an example table using the \texttt{booktabs} package. Note the positioning of the caption above the table rather than below for figures.

\begin{table}[h]
\centering
 \caption{The five oldest Cambridge colleges}
 \label{tab:example-table}
    \begin{tabular}{c c c}
        \toprule
         \textbf{\#} & \textbf{College} & \textbf{Founding date} \\
         \midrule
            1 & Peterhouse & 1284 \\
            2 & Clare & 1326 \\
            3 & Pembroke & 1347 \\
            4 & Gonville and Caius & 1348 \\
            5 & Trinity Hall & 1350 \\
         \bottomrule
    \end{tabular}
\end{table}

To cite, we simply use the \texttt{\textbackslash citep\{\}} or \texttt{\textbackslash citet\{\}} command, in combination with a \textit{bibtex} file.

Kitchen supplies are known to go missing in an office context \citep{CaseOfDisappearingTeaspoons}. However, not all hope is lost, nor should we lose faith in humankind's ability to retain teaspoons. For instance, \citet{TeaspoonSolutionProposal} has proposed an elegant solution which could be considered world-changing.

To list code, use the \href{https://www.overleaf.com/learn/latex/Code_listing#Using_listings_to_highlight_code}{\texttt{lstlisting}} package. This contains all kinds of customisation for how code is linted within your \LaTeX \ source files.

\begin{lstlisting}
import java.util.Random;

public class RandomExample {

    public static void main(String[] args) {
        Random random = new Random();

        int[] numbers = new int[5];

        // Fill array with random numbers
        for (int i = 0; i < numbers.length; i++) {
            numbers[i] = random.nextInt(100);
        }

        // Print the numbers
        System.out.println("Generated numbers:");
        for (int num : numbers) {
            System.out.println(num);
        }

        // Simple calculation
        int sum = 0;
        for (int num : numbers) {
            sum += num;
        }

        double average = (double) sum / numbers.length;
        System.out.println("Average: " + average);
    }
}
\end{lstlisting}